\honest{No One Can Exempt You From Rationality's Laws}{Eliezer Yudkowsky}{2007}

Traditional Rationality states its rules as social obligations. If you want me to accept a belief, you owe me evidence, and a theory owes the world bold predictions and a chance to be falsified. \nb{As in earlier posts, no one who holds this tradition is named, and the falsification rules are Karl Popper's.} This phrasing passes easily by word of mouth, because people detect social cheating far better than breaches of ``isomorphic'' abstract rules. \nb{The finding, from Wason's selection task, is standard. Whether the social and abstract versions are really of the same form is disputed, and the finding is about detecting violations, not about how traditions spread.}

Seen as a social rule, rationality invites a strange defence. Religious people say, ``Well, you can't justify your belief in science!'' In other words: ``You're doing it too!'' \nb{The retort can also be read as a challenge to the standard itself, since science relies on induction. The post does not take up that reading.}

To Bayesians, the brain is an engine that concentrates evidence into a map of the territory, and the rules of rationality are laws like the second law of thermodynamics. A reliable belief needs a calculable amount of evidence, as cooling a refrigerator needs a calculable minimum of energy. A refrigerator could in principle cool itself while generating electricity, and you could draw an accurate street map of New York without visiting, but only with infinitesimally tiny probability. \nb{The probability is tiny because the map is detailed. A guess without evidence about a single yes-or-no question is right about half the time. What it lacks is reliability, which is what the law, as the post first states it, concerns.} Before trying it, wait for a cup of water to freeze by itself.

If the rules were customs, pointing at others might excuse you. They are laws, so it helps as much as listing reasons why you should not fall off a cliff. A judge may excuse two parties who both broke a contract, but two badly designed engines both fail. ``One design error cannot excuse another.'' Even if we all vote that it is unfair for your refrigerator to need electricity, it will not run. Even if we vote that you need not visit New York, the map will be wrong.

The law lets everyone hold their own beliefs, and the wiser countries arrest no one for weird ones. But nature entitles no one to accuracy. Rationalists no more decide the laws of rationality than physicists decide the laws of physics. If you persuaded every physicist that you were exempt from gravity and walked off a cliff, you would fall. Even ``we don't decide'' is too human a phrase: there is no authority that could exempt you. ``There is only cause and effect.'' So when you plead to be excused just this once: ``We can't excuse you. It isn't up to us.''
